\section{Analytic exclusion of nonalternating consecutive tails}
\label{stex:section}

The finite unit-circle theorem can be complemented by an all-order
argument on a range of shifts.  Its classical input is the negative
simple-root property of the Euler--Frobenius numerator on the unit shift
interval; see Janson, Remark~A.2~\cite{stex:Janson}.  We give the short
interlacing proof below, as well as the elementary Enestr\"om--Kakeya
argument needed for a finite shift sum.  The new reduction compares their
angular logarithmic derivatives.  No historical-priority claim is made
for these classical polynomial facts or for the resulting exclusion.

Throughout, $A_k(w,a)=\Phi(w,-k,a)$, $k\ge0$, and $a>0$.
A nonalternating unit-circle color means $|w|=1$, $w\notin\{1,-1\}$.

\subsection{The base interval and a bound for finite shift sums}

\begin{lemma}[Classical negative roots and an angular lower bound]
\label{stex:base}
For $0<b\le1$ and $k\ge1$, the numerator
\[
 F_{k,b}(w)=(1-w)^{k+1}A_k(w,b)
\]
has only simple negative real roots.  Its degree is $k$ when $b<1$
and $k-1$ when $b=1$.  If $w=e^{i\theta}$ with $0<\theta<\pi$, then
\begin{equation}\label{stex:lower}
 \Im\frac{w\partial_w A_k(w,b)}{A_k(w,b)}
       >\frac{k+1}{4}\sin\theta.
\end{equation}
\end{lemma}
\begin{proof}
The numerator recurrence is
\[
 F_{k+1,b}=w(1-w)F'_{k,b}
            +[b+(k+1-b)w]F_{k,b},\qquad F_{0,b}=1.
\]
For $0<b<1$ its constant and leading coefficients are $b^k$ and
$(1-b)^k$.  Suppose the $k$ roots of $F_{k,b}$ are simple and negative,
listed in increasing order.  At each root the next numerator has the
sign of $w(1-w)F'_{k,b}$, which alternates and is opposite to the sign
of $F'_{k,b}$.  Its signs at $-\infty$ and $0$ provide the two remaining
intervals.  Thus $F_{k+1,b}$ has one negative root in each interval
separated by those $k$ roots, accounting for its entire degree.
The induction starts at $F_{1,b}=b+(1-b)w$.
For $b=1$ start instead at $F_{1,1}=1$; the same recurrence has
positive constant and leading coefficients, degree $k-1$ at order $k$,
and the identical sign argument applies.

Write the roots as $-r_j$, $r_j>0$; an empty product is allowed.  Then
\[
 \frac{w\partial_w A_k}{A_k}
   =\sum_j\frac{w}{w+r_j}+(k+1)\frac{w}{1-w}.
\]
For a unit color in the upper half-plane,
\[
 \Im\frac{w}{w+r_j}
       =\frac{r_j\sin\theta}{|w+r_j|^2}>0,
 \qquad
 \Im\frac{w}{1-w}=\frac{\sin\theta}{|1-w|^2}>
                              \frac{\sin\theta}{4}.
\]
This proves \eqref{stex:lower}, including the empty-product case.
\end{proof}

\begin{lemma}[Angular bound for a polynomial whose roots lie in a disc]
\label{stex:upper}
Let $S$ be a real polynomial of degree $d$ whose complex roots have
modulus at most $\rho<1$.  For $w=e^{i\theta}$, $0<\theta<\pi$,
\begin{equation}\label{stex:upper-formula}
 \left|\Im\frac{wS'(w)}{S(w)}\right|
       \le\frac{d\rho}{(1-\rho)^2}\sin\theta.
\end{equation}
In particular, for $n\ge2$, $b>0$, and $k\ge1$, the polynomial
\begin{equation}\label{stex:S}
 S_{k,n,b}(w)=\sum_{j=0}^{n-1}(b+j)^k w^j
\end{equation}
satisfies this bound with
\[
 d=n-1,\qquad
 \rho=\left(\frac{b+n-2}{b+n-1}\right)^k<1.
\]
\end{lemma}
\begin{proof}
If the roots are $z_1,\ldots,z_d$, set $p_m=\sum_jz_j^m$.
The coefficients of $S$ are real, so every $p_m$ is real, and
$|p_m|\le d\rho^m$.  The normally convergent expansion on the unit
circle gives
\[
 \frac{wS'(w)}{S(w)}=d+\sum_{m\ge1}p_m w^{-m}.
\]
Using $|\sin(m\theta)|\le m\sin\theta$ and
$\sum_{m\ge1}m\rho^m=\rho/(1-\rho)^2$ proves the bound.

For completeness, if $S(w)=\sum_{j=0}^d c_jw^j$ with all $c_j>0$,
its roots lie in the disc of radius $\max_{1\le j\le d}c_{j-1}/c_j$.
Scale $w=\rho v$ so that the new coefficients $a_j=c_j\rho^j$
are nondecreasing.  If the scaled polynomial vanished at $|v|>1$,
the identity
\[
 (1-v)\sum_{j=0}^d a_jv^j
    =a_0+\sum_{j=1}^d(a_j-a_{j-1})v^j-a_dv^{d+1}
\]
would give $a_d|v|^{d+1}\le a_d|v|^d$, a contradiction.
For \eqref{stex:S} the successive coefficient ratios increase with $j$;
their maximum is the stated $\rho$.
\end{proof}

The base argument gives a further conclusion outside the unit circle.
For $0<a\le2$, no pair $A_k(w,a),A_{k+1}(w,a)$ can vanish
at any nonreal color $w$. Indeed, in the upper half-plane
$\Im(w/(w+r))=r\Im(w)/|w+r|^2>0$ for $r>0$, and
$\Im(w/(1-w))=\Im(w)/|1-w|^2>0$. For $a\le1$ the
negative-root lemma excludes a single nonreal zero; for
$1<a\le2$ the shift sum is a nonzero constant. The proof
below then applies without a unit-modulus assumption.

\subsection{The exclusion criterion}

\begin{theorem}[A sufficient analytic exclusion criterion]
\label{stex:criterion}
Let $k\ge1$, and write uniquely
\[
 a=b+n,\qquad 0<b\le1,\qquad n\in\mathbb Z_{\ge0}.
\]
There are no simultaneous zeros $A_k(w,a)=A_{k+1}(w,a)=0$ at
nonalternating unit-circle colors if either $n\le1$, or $n\ge2$ and
\begin{equation}\label{stex:criterion-formula}
 (k+1)(1-\rho)^2\ge4(n-1)\rho,
 \qquad
 \rho=\left(\frac{a-2}{a-1}\right)^k.
\end{equation}
For every bounded interval of positive shifts, the criterion excludes
all sufficiently large orders uniformly.
\end{theorem}
\begin{proof}
Because the coefficients are real, complex conjugation reduces the
problem to $w=e^{i\theta}$, $0<\theta<\pi$.  If $n=0$, the
negative-root lemma excludes even a single zero at such a color.
For $n\ge1$, iteration of the Lerch shift identity gives
\begin{equation}\label{stex:shift}
 w^nA_k(w,b+n)=A_k(w,b)-S_{k,n,b}(w).
\end{equation}
Moreover
\[
 w\partial_w A_k(w,a)=A_{k+1}(w,a)-aA_k(w,a).
\]
Consequently simultaneous zeros imply both
$A_k(w,b)=S_{k,n,b}(w)$ and
$w\partial_w A_k(w,b)=wS'_{k,n,b}(w)$.
The common value is nonzero by Lemma~\ref{stex:base}, so their
logarithmic derivatives are equal.  If $n=1$, $S$ is the nonzero
constant $b^k$, contradicting \eqref{stex:lower}.
For $n\ge2$, Lemma~\ref{stex:upper} and
\eqref{stex:criterion-formula} contradict the strict lower bound
\eqref{stex:lower}.  Notice that equality in the displayed criterion
still excludes a zero.

For fixed $a>2$, the ratio $(a-2)/(a-1)$ is strictly between zero and
one.  On a bounded shift interval it is bounded above by a constant
less than one, and $n$ is bounded.  Thus $\rho$ decays geometrically
with $k$, while $k+1$ increases, proving the final assertion.
\end{proof}

\begin{corollary}[All-order exclusion on $0<a\le7$]
\label{stex:seven}
For every $k\ge0$, $0<a\le7$, and $|w|=1$ with
$w\notin\{1,-1\}$, the pair $A_k(w,a),A_{k+1}(w,a)$ cannot
vanish simultaneously.
\end{corollary}
\begin{proof}
The case $k=0$ follows from $A_0=(1-w)^{-1}\ne0$.
For $1\le k\le7$, apply the published unit-circle classification
through order eight, Theorem~\ref{stcol:circle-eight}, with $q=k+1$.
For $k\ge8$, write $a=b+n$ as above.  Then $n\le6$.
The cases $n\le1$ are already covered.  Otherwise
$n-1\le5$ and $\rho\le(5/6)^k\le(5/6)^8$.
The exact rational margin is
\[
 9\left(1-\left(\frac56\right)^8\right)^2
       -20\left(\frac56\right)^8
   =\frac{1831480182729}{2821109907456}>0.
\]
As $k$ increases, the left positive term in
$(k+1)(1-\rho)^2-4(n-1)\rho$ increases and the subtracted term
decreases.  The sufficient criterion therefore holds for all $k\ge8$.
The all-order conclusion depends on this monotonic argument and on
the previously certified finite range; it is not an extrapolation
of a finite search.
\end{proof}

\begin{corollary}[A simple linear shift-order region]
\label{stex:linear}
For $k\ge1$, nonalternating unit-circle consecutive zeros are impossible
whenever $0<a\le1+2k/3$.
\end{corollary}
\begin{proof}
Only $a>2$ needs consideration.  Then $n-1<a-1\le2k/3$, and
\[
 \rho=\left(1-\frac1{a-1}\right)^k
       <e^{-k/(a-1)}\le e^{-3/2}<\frac9{40}.
\]
The last bound follows, for example, by comparing the first seven
positive Taylor terms of $e^{3/2}$ with $40/9$.
Since $u/(1-u)^2$ increases on $(0,1)$,
\[
 \frac{4(n-1)\rho}{(1-\rho)^2}
       <\frac{8k}{3}\frac{9/40}{(1-9/40)^2}
       =\frac{960}{961}k<k+1.
\]
The criterion applies.
\end{proof}

On the strip $0<a\le7$, every fixed-shift homogeneous balance at a
nonalternating unit-circle color is therefore unique, and each colored
kernel has at most one extra ordinary positive-integer Mellin moment.
The algebraic order-fifteen point, with $a_*\approx5.94585$, lies in
this strip and has precisely one such moment, so the one-moment bound
is attained.  The theorem excludes consecutive vanishing, not single
vanishing.  At the alternating color the already established Brillhart
classification still supplies the golden-shift exception.
Outside the sufficient regions above, the global nonalternating
consecutive-tail problem remains open in this investigation.
